% GENERATED FILE. Edit canonical lessons, then run npm run generate:books.
% canonical-source-hash: fnv1a64:2321443ca2c6c19c
% canonical-lessons: GE-C384-genehmigung, GE-C384-genehmigen, GE-C384-gueltig, GE-C384-verlaengern, GE-C384-ablaufen

\chapter{Die Genehmigung}
\label{ch:ge-die-genehmigung}
\hlchaptermodality{384}

\begin{chapteropening}
\textbf{By the end of this chapter:} \emph{I can talk about permits: getting approval, how long a document is valid, renewing it and when it expires.}

Complete the last lesson of chapter 384: I can talk about permits: getting approval, how long a document is valid, renewing it and when it expires.
\end{chapteropening}

\section[die Genehmigung]{die Genehmigung — the permit, the approval}

\label{lesson:GE-C384-genehmigung}



\begin{quote}
Before the new one: say the German for the form, then the German for the signature.
\end{quote}

\subsection*{You'll want to know: die Genehmigung}
\textbf{die Genehmigung} — \textquotedblleft{}the permit, the approval\textquotedblright{}.

\begin{quote}
Für das Straßenfest braucht man eine Genehmigung der Stadt.
\end{quote}

\begin{quote}
\emph{For the street party you need a permit from the town.}
\end{quote}

Plural: \emph{die Genehmigungen}.

\begin{cousinweb}
\emph{Genehmigung} is the noun of \emph{genehmigen} \textquotedblleft{}to approve\textquotedblright{}, built on the old adjective \emph{genehm} \textquotedblleft{}acceptable\textquotedblright{}, of the family of \emph{nehmen} \textquotedblleft{}to take\textquotedblright{}, Proto-Germanic \emph{nemaną} (EtymDB): something you can take.
\end{cousinweb}

\begin{grammarlens}[title={What you've built}]
One more word for telling, explaining and describing at B1.
\end{grammarlens}

\subsection*{Your turn}
\begin{itemize}
\raggedright
  \item \emph{Say it:} \emph{die Genehmigung}
  \item \emph{Say it:} \emph{die Genehmigung}, once more
  \item \emph{Say it:} a sentence of your own with \emph{die Genehmigung}
  \item \emph{Recall:} say the German for the form, then the German for the signature, then say \emph{die Genehmigung} again
\end{itemize}

\begin{tcolorbox}[breakable,colback=teal!4,colframe=teal!35!black,title={Before you move on}]
What does die Genehmigung mean? (\textquotedblleft{}The permit, the approval\textquotedblright{}.) Say it once more.
\end{tcolorbox}
\section[genehmigen]{genehmigen — to approve, to permit}

\label{lesson:GE-C384-genehmigen}



\begin{quote}
Before the new one: say the German for the signature, then the German for the permit.
\end{quote}

\subsection*{You'll want to know: genehmigen}
\textbf{genehmigen} — \textquotedblleft{}to approve, to permit\textquotedblright{}.

\begin{quote}
Mein Chef hat meinen Urlaub genehmigt.
\end{quote}

\begin{quote}
\emph{My boss has approved my holiday.}
\end{quote}

Inseparable: \emph{er genehmigt, er hat genehmigt}. \emph{Sich etwas genehmigen} — to treat yourself to something.

\begin{cousinweb}
\emph{genehmigen} is built on the old adjective \emph{genehm} \textquotedblleft{}acceptable\textquotedblright{}, of the family of \emph{nehmen} \textquotedblleft{}to take\textquotedblright{}, Old High German \emph{neman}, Proto-Germanic \emph{nemaną} (EtymDB).
\end{cousinweb}

\begin{grammarlens}[title={What you've built}]
One more word for telling, explaining and describing at B1.
\end{grammarlens}

\subsection*{Your turn}
\begin{itemize}
\raggedright
  \item \emph{Say it:} \emph{genehmigen}
  \item \emph{Say it:} \emph{genehmigen}, once more
  \item \emph{Say it:} a sentence of your own with \emph{genehmigen}
  \item \emph{Recall:} say the German for the signature, then the German for the permit, then say \emph{genehmigen} again
\end{itemize}

\begin{tcolorbox}[breakable,colback=teal!4,colframe=teal!35!black,title={Before you move on}]
What does genehmigen mean? (\textquotedblleft{}To approve, to permit\textquotedblright{}.) Say it once more.
\end{tcolorbox}
\section[gültig]{gültig — valid}

\label{lesson:GE-C384-gueltig}



\begin{quote}
Before the new one: say the German for the permit, then the German for to approve.
\end{quote}

\subsection*{You'll want to know: gültig}
\textbf{gültig} — \textquotedblleft{}valid\textquotedblright{}.

\begin{quote}
Die Fahrkarte ist nur heute gültig.
\end{quote}

\begin{quote}
\emph{The ticket is only valid today.}
\end{quote}

The opposite is \emph{ungültig} — invalid.

\begin{cousinweb}
\emph{gültig} is inherited from Middle High German \emph{gültic} (EtymDB), from \emph{gülte} \textquotedblleft{}income, payment\textquotedblright{}, of the family of \emph{gelten} \textquotedblleft{}to be worth\textquotedblright{}, Old High German \emph{geltan} \textquotedblleft{}to pay\textquotedblright{}.
\end{cousinweb}

\begin{grammarlens}[title={What you've built}]
One more word for telling, explaining and describing at B1.
\end{grammarlens}

\subsection*{Your turn}
\begin{itemize}
\raggedright
  \item \emph{Say it:} \emph{gültig}
  \item \emph{Say it:} \emph{gültig}, once more
  \item \emph{Say it:} a sentence of your own with \emph{gültig}
  \item \emph{Recall:} say the German for the permit, then the German for to approve, then say \emph{gültig} again
\end{itemize}

\begin{tcolorbox}[breakable,colback=teal!4,colframe=teal!35!black,title={Before you move on}]
What does gültig mean? (\textquotedblleft{}Valid\textquotedblright{}.) Say it once more.
\end{tcolorbox}
\section[verlängern]{verlängern — to extend, to renew}

\label{lesson:GE-C384-verlaengern}



\begin{quote}
Before the new one: say the German for to approve, then the German for valid.
\end{quote}

\subsection*{You'll want to know: verlängern}
\textbf{verlängern} — \textquotedblleft{}to extend, to renew\textquotedblright{}.

\begin{quote}
Ich muss meinen Ausweis verlängern lassen.
\end{quote}

\begin{quote}
\emph{I have to have my identity card renewed.}
\end{quote}

Inseparable: \emph{er verlängert, er hat verlängert}. Built on \emph{länger} — longer.

\begin{cousinweb}
\emph{verlängern} is built with \emph{ver-} on \emph{länger} \textquotedblleft{}longer\textquotedblright{}, from \emph{lang} \textquotedblleft{}long\textquotedblright{}, inherited from Old High German \emph{lang} and Proto-Germanic \emph{langaz} (EtymDB): you make the time longer.
\end{cousinweb}

\begin{grammarlens}[title={What you've built}]
One more word for telling, explaining and describing at B1.
\end{grammarlens}

\subsection*{Your turn}
\begin{itemize}
\raggedright
  \item \emph{Say it:} \emph{verlängern}
  \item \emph{Say it:} \emph{verlängern}, once more
  \item \emph{Say it:} a sentence of your own with \emph{verlängern}
  \item \emph{Recall:} say the German for to approve, then the German for valid, then say \emph{verlängern} again
\end{itemize}

\begin{tcolorbox}[breakable,colback=teal!4,colframe=teal!35!black,title={Before you move on}]
What does verlängern mean? (\textquotedblleft{}To extend, to renew\textquotedblright{}.) Say it once more.
\end{tcolorbox}
\section[ablaufen]{ablaufen — to expire; to go, to proceed}

\label{lesson:GE-C384-ablaufen}



\begin{quote}
Before the new one: say the German for valid, then the German for to extend.
\end{quote}

\subsection*{You'll want to know: ablaufen}
\textbf{ablaufen} — \textquotedblleft{}to expire; to go, to proceed\textquotedblright{}.

\begin{quote}
Mein Ausweis läuft im Mai ab.
\end{quote}

\begin{quote}
\emph{My identity card expires in May.}
\end{quote}

Separable, Perfekt with \emph{sein}: \emph{er läuft ab}, \emph{er ist abgelaufen}. Also to go: \emph{Wie ist das Gespräch abgelaufen?} — how did the conversation go?

\begin{cousinweb}
\emph{ablaufen} is built from \emph{ab} \textquotedblleft{}off, away\textquotedblright{} and \emph{laufen} \textquotedblleft{}to run\textquotedblright{}, inherited from Old High German \emph{hloufan} and Proto-Germanic \emph{hlaupaną} (EtymDB): the time runs out.
\end{cousinweb}

\begin{grammarlens}[title={What you've built}]
One more word for telling, explaining and describing at B1.
\end{grammarlens}

\subsection*{Your turn}
\begin{itemize}
\raggedright
  \item \emph{Say it:} \emph{ablaufen}
  \item \emph{Say it:} \emph{ablaufen}, once more
  \item \emph{Say it:} a sentence of your own with \emph{ablaufen}
  \item \emph{Recall:} say the German for valid, then the German for to extend, then say \emph{ablaufen} again
\end{itemize}

\begin{tcolorbox}[breakable,colback=teal!4,colframe=teal!35!black,title={Before you move on}]
What does ablaufen mean? (\textquotedblleft{}To expire; to go, to proceed\textquotedblright{}.) Say it once more.
\end{tcolorbox}
